\documentclass[10pt,a4paper]{article}
\usepackage[margin=2.2cm]{geometry}
\usepackage{graphicx}
\usepackage{booktabs}
\usepackage{siunitx}
\usepackage{amsmath}
\usepackage{float}
\usepackage{caption}
\usepackage[hidelinks]{hyperref}
\sisetup{detect-all, table-number-alignment=center}
\DeclareSIUnit{\dBc}{dBc}
\DeclareSIUnit{\dBFS}{dBFS}
\setlength{\emergencystretch}{2em}

\input{results.tex}

\title{SARADC (IHP SG13G2): transistor-level characterization by simulation}
\author{saradc\_auto}
\date{\today\\[2pt]\small Data set: \texttt{\CharDir}}

\begin{document}
\maketitle

\section*{Summary}
\begin{center}
\begin{tabular}{ll}
\toprule
Resolution & 5 bits \\
Clock / sample rate & \SI{\Fclk}{MHz} / \SI{\Fs}{MS/s} (7 clocks per conversion) \\
Simulator, corner & \Simulator, \texttt{\Corner}, \SI{25}{\celsius}, schematic netlist (post-layout: Section~\ref{sec:postlayout}) \\
PDK & \PdkName, git \texttt{\PdkCommit} (\PdkBranch, \PdkDate) \\
Supply, references, input CM & \SI{\Vdd}{V}; \SI{\Vrefl}{V} / \SI{\Vrefh}{V}; \SI{\Vcm}{V} (as on the test board) \\
LSB (end point) & \SI{\Lsb}{mV} differential (ideal \SI{\LsbIdeal}{mV}, gain error \SI{\GainErr}{\percent}) \\
Offset & \num{\OffsetLsb} LSB \\
DNL & \numrange{\DnlMin}{\DnlMax} LSB \\
INL (end point) & \numrange{\InlMin}{\InlMax} LSB \\
Missing codes & \MissingCodes \\
\IfFileExists{table_sweep.tex}{ENOB at \SI{\SineLowFreq}{kHz} / \SI{\SineHighFreq}{kHz} & \SineLowEnob{} / \SineHighEnob{} bits \\}%
{ENOB & sine records still running (see Section~\ref{sec:dynamic}) \\}
\bottomrule
\end{tabular}
\end{center}

\section{Circuit}
The converter is a fully differential, charge-redistribution SAR of 5 bits assembled from standard cells
(\texttt{rtl/saradc\_analog.v}). Each side has an MSB capacitive DAC that samples the input
(\texttt{SARADC\_MSB\_CDAC}, 16 units, connected backwards so that the input is sampled on the
plates that later move) and a binary LSB DAC of 1, 2, 4 and 8 units (plus one dummy unit) that switches
between $V_{REFH}$ and $V_{REFL}$. A unit is six \texttt{SARADC\_CELL\_INVX16\_ASCAP} capacitors: an
inverter whose drain and both source terminals are tied together (terminal FL) and whose gate is the
switched plate (terminal VSH), so the capacitance is a MOS gate capacitance in series with the junctions of
FL to the wells. The switches are transmission gates of \texttt{SARADC\_CELL\_INVX0\_ASSW}. The dynamic
comparator compares the two DAC outputs, and the digital part (\texttt{sar\_logic\_wreset}) takes
7 clocks per conversion: one sampling clock (SAMPLE high), five bit decisions and one clock with VALID high.

\section{Method}
\begin{itemize}
  \item \textbf{Conditions}: those of the chip measurement (adapter board
        \texttt{IHP130\_ASICONE\_ADP\_\allowbreak QFN48\_0p5Q\_Adapter}): every supply pin of the macro at
        $V_{DD} = \SI{\Vdd}{V}$, $V_{REFH} = \SI{\Vrefh}{V}$, $V_{REFL} = \SI{\Vrefl}{V}$, input common mode
        \SI{\Vcm}{V}; differential full scale $2(V_{REFH}-V_{REFL})$ peak-to-peak, so the ideal LSB is
        \SI{\LsbIdeal}{mV}. Clock \SI{\Fclk}{MHz} as on the chip measurement, the inputs are driven
        through \SI{50}{\ohm}. The testbench of the design (\texttt{src/saradc\_tb\_body.sp}) uses
        \SI{1.8}{V}, \SI{0.45}{V}/\SI{1.35}{V} and \SI{0.9}{V}; the static transfer was also simulated there
        (Section~\ref{sec:static}). The options \texttt{--vdd}, \texttt{--vrefh}, \texttt{--vrefl} and
        \texttt{--vcm} of \texttt{adc\_char.py} select the conditions.
  \item \textbf{Simulator}: ngspice-47~\cite{ngspice} with the PSP103 OSDI model of the PDK~\cite{ihppdk} (compiled with
        OpenVAF), KLU, Gear integration, \texttt{reltol=1e-4}, \texttt{trtol=7}. With the default
        \texttt{reltol} the timestep collapses at some clock edges; the tighter tolerance runs through
        and is retried at \texttt{1e-5} if needed. Xyce (same PSP103 model through its ADMS plugin, KLU,
        one process per job) converts the same inputs of a 3-conversion test job to the same codes,
        except exactly at $v_d = 0$, which lies within \SI{0.5}{mV} of the transition $T_{16}$, but
        takes about three times longer per conversion (about ten times with its default iterative
        solver), so the campaign uses ngspice.
  \item \textbf{Netlist}: the post-route schematic netlist \texttt{pnr/outputs/SARADC.cdl}
        (no layout parasitics), flattened. Transistors that are exactly in parallel (the 12+12 fingers
        of each unit capacitor and the switches of one DAC unit, which have the same nodes, model and
        size) are merged into one instance with the summed multiplier (\texttt{merge\_parallel} in
        \texttt{src/adc\_char.py}). This is exact for the PSP model and reduces the circuit from
        24\,198 to 2\,902 transistors; the KLayout netlist comparer confirms that the merged netlist
        is the same circuit as the schematic.
  \item \textbf{Testbench}: the capture registers and the ideal reconstruction DAC of the original
        testbench only observe RESULT; they are replaced by their load (\SI{4}{fF} per line) and the
        codes are read directly on RESULT in the middle of each VALID pulse.
  \item \textbf{Jobs}: the records are split into short independent transient runs (one process each,
        about 30 in parallel). Each job simulates the reset, one warm-up conversion and $K$ conversions.
        For the static transfer the input is a DC staircase (one level per conversion, changed in the
        middle of the previous conversion); for the dynamic tests it is a slice of a coherent sine whose
        phase is shifted so that the job converts the right samples of the record.
  \item \textbf{Static transfer}: a staircase of \SI{0.5}{LSB} steps over the full scale, then five
        bisection rounds per transition: the transition levels $T_k$ are known to
        \SI{\Resolution}{mV}. DNL and INL use the end-point line through $T_1$ and $T_{31}$ (IEEE~1241~\cite{ieee1241}), as the
        measurement.
  \item \textbf{Dynamic}: records of \IfFileExists{table_sweep.tex}{\NSine{}}{256} samples with an odd prime number of cycles, amplitude
        \IfFileExists{table_sweep.tex}{\SineAmpLsb{}}{13.5} LSB (peak) around mid scale, close to the amplitude of the chip measurement.
        SINAD from a sine fit, harmonics from the FFT (coherent, rectangular window).
        SINAD$_{st}$ / ENOB$_{st}$ is what the simulated static transfer alone predicts for the same
        input samples.
  \item \textbf{Board front end}: the single-ended to differential converter of the test board
        (TI SBOA216~\cite{sboa216} with LMV358 op-amps~\cite{lmv358}, \SI{47.5}{\ohm} per side,
        \SI{1.2}{nF} across the pair and \SI{120}{pF} to ground) is not simulated at transistor
        level. Its linear model (single-pole op-amps, GBW \SI{1.4}{MHz}, fitted to the measured
        amplitude response; \texttt{--frontend}) gives the amplitude and phase of each input,
        which the \texttt{sinefe} records apply to the ADC \IfFileExists{table_fe.tex}{(Section~\ref{sec:fe})}{(Section~\ref{sec:dynamic})}.
  \item \textbf{Not included}: thermal and flicker noise (transient simulation), wire resistance,
        mismatch, the reference drivers of the board, and the nonlinearity and slew rate of
        the op-amps of the front end.
  \item \textbf{PDK}: IHP-Open-PDK~\cite{ihppdk}, branch \texttt{\PdkBranch}, commit \texttt{\PdkCommit}
        of \PdkDate{} (\texttt{git describe} gives \PdkDescribe). The compiled device models (PSP103
        and r3\_cmc as OSDI for ngspice, as plugins for Xyce) were rebuilt from the Verilog-A sources of
        this commit.
  \item \textbf{Layout parasitics}: not in the schematic simulation. The capacitances of the routed
        layout, obtained by a split extraction (Section~\ref{sec:split}), are simulated separately
        (static transfer only, Section~\ref{sec:postlayout}).
\end{itemize}

\section{Unit capacitor}
The unit capacitance seen from the gate (from the AC runs of \texttt{make test\_cap\_linearity}) changes with
the bias between \SI{\CapMin}{fF} (at \SI{\CapVmin}{V}) and \SI{\CapMax}{fF}. Because the MSB DAC
samples the input on these capacitors, the charge sampled is a nonlinear function of the input.
\begin{figure}[H]\centering
\includegraphics{fig/cap_cv.pdf}
\caption{Capacitance of \texttt{SARADC\_CELL\_INVX16\_ASCAP} versus the bias of the gate, at \SI{100}{MHz}.}
\end{figure}

\section{Static linearity}
\label{sec:static}
\begin{figure}[H]\centering
\includegraphics{fig/transfer.pdf}
\caption{Simulated transfer characteristic (transition levels from the bisection).}
\end{figure}
\begin{figure}[H]\centering
\includegraphics{fig/dnl_inl.pdf}
\caption{DNL and INL (end point).}
\end{figure}
\begin{table}[H]\centering\small
\input{table_codes.tex}
\caption{Transition levels (differential input) and linearity per code.}
\end{table}

\section{Dynamic performance}
\label{sec:dynamic}
\IfFileExists{table_sweep.tex}{%
\begin{figure}[H]\centering
\includegraphics{fig/sine_time.pdf}
\caption{Output for the \SI{\SineLowFreq}{kHz} sine (first three periods).}
\end{figure}
\begin{figure}[H]\centering
\includegraphics{fig/spectrum_low.pdf}
\caption{Spectrum at \SI{\SineLowFreq}{kHz}.}
\end{figure}
\begin{figure}[H]\centering
\includegraphics{fig/spectrum_high.pdf}
\caption{Spectrum at \SI{\SineHighFreq}{kHz}.}
\end{figure}
\begin{figure}[H]\centering
\includegraphics{fig/sweep_enob.pdf}
\caption{ENOB versus input frequency.}
\end{figure}
\begin{figure}[H]\centering
\includegraphics{fig/sweep_harmonics.pdf}
\caption{THD and harmonics versus input frequency. H2 is below \SI{-100}{dBc} (a numerical zero) at all
frequencies and is not drawn.}
\end{figure}
\begin{table}[H]\centering\scriptsize
\input{table_sweep.tex}
\caption{Dynamic results. $A$ is the fitted output amplitude.}
\end{table}}{%
The sine records (six frequencies with the ideal differential input and four through the model of the
board front end, 256 samples each) are still running at the time of this version of the report;
this section is filled in when they finish.}

\IfFileExists{table_fe.tex}{%
\section{Input through the board front end}
\label{sec:fe}
The same sines, with each input given the amplitude and phase of the linear model of the board front
end. The model keeps the differential signal a pure sine but makes the two inputs no longer
opposite: the input common mode moves at the signal frequency (a few \si{\percent} of the
differential amplitude at \SI{17}{kHz}, about half at \SI{630}{kHz}).
\begin{table}[H]\centering\small
\input{table_fe.tex}
\caption{Ideal differential input against the input through the front-end model.}
\end{table}}{}

\section{Layout parasitic extraction}
\label{sec:pex}
The post-layout simulation needs a netlist of the routed layout (\texttt{signoff/outputs/SARADC.gds},
\SI{566}{\micro m} $\times$ \SI{361}{\micro m}, placed flat: 24\,198 transistors in standard cells and
in the 720 unit-capacitor and 1\,452 switch cells, plus about 30\,000 via instances). Two extractors were
tried with the PDK above; neither gave a netlist that can be simulated.

\subsection{Magic and Netgen}
Magic~\cite{magic} 8.3.623 with the technology file of the PDK (\texttt{libs.tech/magic}) extracts the
layout (\texttt{pex/outputs/magic}) and finds the same devices as the schematic (12\,093 nMOS and 12\,105
pMOS). The result is not usable:
\begin{itemize}
  \item \textbf{Capacitances.} The 71\,863 extracted capacitors add up to \SI{4.27}{nF} on AVDD,
        \SI{1.04}{nF} on VSS, \SI{516}{pF} and \SI{423}{pF} on the inputs VIN and VIP and about
        \SI{415}{pF} on the output of a single switch cell (\texttt{SARADC\_CELL\_INVX0\_ASSW}). The
        whole capacitive DAC is of the order of \SI{10}{pF} (720 units of about \SI{11}{fF}), so these
        values are two orders of magnitude too large. This points to the capacitance rules of the
        magic technology file for the current layer set of the PDK; it was not investigated further.
  \item \textbf{LVS with Netgen.} Netgen~\cite{netgen} 1.5.287 cannot compare the extracted netlist with
        the schematic: magic writes the transistors as calls to the subcircuits
        \texttt{sg13\_lv\_nmos}/\texttt{sg13\_lv\_pmos}, which are not defined in the extracted netlist,
        so Netgen gives them generic pins 1 to 4 while the schematic has drain, gate, source and bulk.
        This is a naming mismatch between the two netlist formats, not a connectivity error; the KLayout
        LVS of the same GDS passes.
  \item \textbf{Net names.} Internal nets are named after cell pins (e.g.\
        \texttt{SARADC\_CELL\_INVX0\_ASSW\_99/zn}); the nets of the DAC (the top plates VOUTH and
        VOUTL, the plates of each bit) cannot be identified, so not even the coupling capacitances of the
        DAC can be read from it.
\end{itemize}

\subsection{KLayout-PEX}
KLayout-PEX~\cite{kpex} 0.3.12 on KLayout~\cite{klayout} 0.30.10 (\texttt{pex/outputs/kpex\_full}) first runs the
LVS of the PDK:
\begin{itemize}
  \item \textbf{LVS.} It fails with the original schematic because the layouts of the
        \texttt{sg13g2f} standard cells have no pin labels. The analog cells (\texttt{INVX0\_ASSW},
        \texttt{INVX16\_ASCAP}) match as they are. With the standard cells flattened in a copy of the
        schematic (\texttt{pex/outputs/kpex\_full/SARADC.flatstd.spice}) the netlists match.
  \item \textbf{2.5D capacitance extraction} (\texttt{--2.5D --mode CC}). The engine works on the
        flat layout on a single thread and wrote nothing within the time limits:
        \begin{center}\small
        \begin{tabular}{lll}
        \toprule
        Run & Options & Result \\
        \midrule
        1 & default & stopped at the \SI{4}{h} limit \\
        2 & default, from the clean LVS database & stopped at the \SI{12}{h} limit \\
        3 & \texttt{--halo 1} (\SI{1}{\micro m}) & stopped at the \SI{10}{h} limit \\
        4 & \texttt{--halo 1}, filler cells without metal removed & stopped by hand after \SI{16}{h}, still in the 2.5D engine \\
        \bottomrule
        \end{tabular}
        \end{center}
        In run 4 the 6\,530 \texttt{SARADC\_FILL1\_NOPOWER} cells, which contain only well and implant
        shapes, are removed from a copy of the layout (\texttt{pex/outputs/kpex\_full/SARADC.nofill.gds}).
  \item The field-solver engines (FasterCap, FastCap2) were not tried: they are much slower than the
        2.5D engine at this size.
\end{itemize}
The commands of all the runs above (magic, netgen and kpex) are in \texttt{pex/} of
\texttt{saradc\_auto} (\texttt{pex/README.md}); the simulation scripts are in \texttt{sim/}
(\texttt{sim/README.md}).

\subsection{Split extraction}
\label{sec:split}
KLayout-PEX does extract small layouts: a window of one tile of the DAC takes \SI{25}{s}. Its time
grows faster than the area (a single window of \SI{895}{\micro m^2}: \SI{25}{s}, \SI{1.2}{GB};
\SI{3582}{\micro m^2}: \SI{141}{s}; \SI{10746}{\micro m^2}: \SI{742}{s}; \SI{30919}{\micro m^2}, half
of a DAC array: \SI{3691}{s} and \SI{30}{GB}, about area$^{1.5}$), so the layout is extracted as many
small windows in parallel and the window netlists are joined again (\texttt{pex/split.py},
\texttt{sim/pex/stitch.py}):
\begin{enumerate}
  \item The GDS is clipped into a grid of windows. The cut lines only cross fillers and taps: between the
        tiles of the DAC (pitch \SI{42.16}{\micro m} $\times$ \SI{21.24}{\micro m}, so no unit capacitor is
        cut) and on row boundaries in the channel and in the digital part.
  \item Every wire segment of the DEF inside a window, and every via landing on a cell pin, gets a text
        label with its net name, so that the part of a net inside a window keeps its name even when it
        only connects outside the window.
  \item The KLayout LVS deck of kpex is run on each window with its unconditional \texttt{purge}
        removed (it drops nets without devices, e.g.\ a stripe crossing a window, which kpex then never
        extracts), then \texttt{kpex --lvsdb --2.5D --mode CC} extracts the coupling and ground
        capacitances (no resistance).
  \item The window netlists are joined by net name: a net cut by a window border is the same net,
        the internal nets of a cell stay local to its window, the substrate is VSS. The joined netlist is
        checked against the flattened schematic: the same transistors, and on every named net the same
        total width on gates, sources/drains and bulks.
\end{enumerate}
Two layouts are extracted:
\begin{itemize}
  \item \textbf{Tapeout} (\texttt{signoff/outputs.tapeout/SARADC.gds}, the fabricated layout): 81
        windows, \SI{6.6}{min} with 16 jobs. 24\,198 transistors and 68\,519 capacitors; the check against
        the schematic passes on all 572 named nets.
  \item \textbf{Hierarchical re-layout} (\texttt{pnr/}): the same circuit placed again with one half of
        the DAC as a macro, used twice (the second half a mirror of the first). The macro (30 windows,
        \SI{2.3}{min}) and the top level without the macros (60 windows, \SI{2.4}{min}) are extracted
        separately and joined; 24\,208 transistors (10 more tie cells) and 57\,911 capacitors, check passed on
        all 569 named nets. Not extracted: the coupling between top-level wires and the inside of the
        macros (only the pin stubs and one supply rail overlap them).
\end{itemize}
In both, kpex on the whole layout does not finish (the whole macro alone: stopped after \SI{4}{h}).

\IfFileExists{results_pex.tex}{% Post-layout static results. Macros, tables and figure from make_sim_report.py --pex
\input{results_pex.tex}
\section{Post-layout simulation}
\label{sec:postlayout}
The joined netlists of Section~\ref{sec:split} (transistors and coupling capacitances, no resistance)
replace the schematic netlist in the same static characterization: same conditions as in the
summary, staircase and bisection of all 31 transition levels to
\SI{0.39}{mV}. With 68\,519 capacitors on top of the transistors one conversion takes several minutes
instead of a few seconds, so only the static transfer is simulated (taped-out netlist: \SI{7.2}{h}
with 26 jobs). The
sine records are not: at low frequency the dynamic part adds nothing (Section~\ref{sec:dynamic}), and the
ENOB of a sine quantized with the post-layout transition levels (ENOB$_{st}$: a coherent sine of 27 LSB
peak to peak, in the converter's own end-point LSB) is what a low-frequency record would give.

\begin{table}[H]\centering\small
\input{table_pex.tex}
\caption{Static results of the schematic and of the post-layout netlists: the bit weights of a
least-squares fit to the transition levels, in the end-point LSB, the end-point LSB, the INL and the
ENOB of a sine quantized with the transition levels. The last two rows are the taped-out netlist with
only the capacitors that touch a switched or floating plate of the DAC (VSH, FL of every bit), and with
all the other capacitors.}
\label{tab:pex}
\end{table}

\begin{figure}[H]\centering
\includegraphics{fig/pex_dnl_inl.pdf}
\caption{DNL and INL of the schematic and of the two post-layout netlists, each in its own end-point
LSB.}
\label{fig:pex_dnl_inl}
\end{figure}

\subsection{Result}
\begin{itemize}
  \item \textbf{The layout breaks the binary weights.} The taped-out netlist has bit weight errors of
        \TpBitEA{}/\TpBitEB{}/\TpBitEC{}/\TpBitED{}/\TpBitEE{}\,\% (bit 0 to 4), DNL from
        \num{\TpDnlMin} to \num{\TpDnlMax} LSB and \TpInlRms{} LSB rms of INL. The ENOB$_{st}$ falls from
        \SchStEnob{} bits (schematic) to \TpStEnob{} bits. The hierarchical re-layout gives
        \HiBitEA{}/\HiBitEB{}/\HiBitEC{}/\HiBitED{}/\HiBitEE{}\,\%, with the same signs; its INL is smaller
        (\HiInlRms{} LSB rms, ENOB$_{st}$ \HiStEnob{} bits) because its two halves are exact mirrors.
  \item \textbf{The DAC is attenuated.} The end-point LSB grows from \SI{\Lsb}{mV} to \SI{\TpLsb}{mV}
        (taped-out, gain error \SI{\TpGainErr}{\percent}) and \SI{\HiLsb}{mV} (re-layout): the wiring
        capacitance of the top plates (VOUTH/VOUTL to the supplies and to each other) does not take part in
        the redistribution and divides the DAC step.
  \item \textbf{Pattern.} Bit 0 is the most reduced, bits 1 and 3 are smaller and bits 2 and 4 larger
        than binary. The pattern follows the plate wiring (Table~\ref{tab:par_tapeout}): bits 1 and 3 have
        the extra bridges on Metal5/TopMetal1 and about \SI{10}{\percent} more wiring capacitance per unit
        on both plates; the MSB has its switched plate (VSH) coupled \SI{11}{fF} per unit to the top plate
        of its own half (VOUTH for H, VOUTL for L), against about \SI{5}{fF} to each top plate for the other bits; the floating plates
        (FL) of bits 0 to 3 couple 7 to \SI{10}{fF} per unit to the top plate of the other half, the MSB
        less than \SI{1}{fF}. These couplings are of the order of \SI{10}{\percent} of the unit
        capacitance (\SI{67.8}{fF} with FL floating), so they change the effective charge of each bit by
        several percent, differently for each bit.
  \item \textbf{Decomposition.} Keeping only the capacitors on the plates (VSH, FL) gives
        \PlBitEA{}/\PlBitEB{}/\PlBitEC{}/\PlBitED{}/\PlBitEE{}\,\% (INL \PlInlRms{} LSB rms);
        \ifdefined\NpBitEA
        keeping all the others gives \NpBitEA{}/\NpBitEB{}/\NpBitEC{}/\NpBitED{}/\NpBitEE{}\,\%
        (INL \NpInlRms{} LSB rms).
        \else
        the run with all the other capacitors is not finished.
        \fi
        Each part alone already has the pattern of the full netlist; the effects do not simply add
        (the sum of the two is larger than the errors of the full netlist). The attenuation of the DAC comes only from the
        plates (LSB \SI{\PlLsb}{mV} with only the plate capacitors, \SI{\NpLsb}{mV} without them).
        Without the plate capacitors the errors come from the other parasitics around the DAC: the
        capacitance of the top plates and of the switch control nets, which couple the switching of
        each bit into the comparator inputs.
\end{itemize}

\begin{table}[H]\centering\scriptsize
\input{table_par_tapeout.tex}
\caption{Wiring capacitance of the plates of the taped-out layout per unit capacitor (\si{fF}), split
by the net on the other side (\texttt{sim/pex/dacpar.py}). VSH is the switched plate of a bit (to the
input, VOUTH or VOUTL), FL the floating plate. VOUTH and VOUTL are the two comparator inputs (top
plates); H is the lower half, whose MSB connects to VOUTH, L the upper half (MSB on VOUTL, the other bits
with VOUTH and VOUTL exchanged). ``VSH--FL'' is the wiring coupling between the two plates of the same bit, in parallel with the
unit capacitors.}
\label{tab:par_tapeout}
\end{table}
\IfFileExists{table_par_hier.tex}{%
In the hierarchical re-layout the L half is an exact mirror of H (the same numbers with VOUTH and VOUTL
exchanged, Table~\ref{tab:par_hier}). In the taped-out layout it is not: the floating plates of the L
half couple about \SI{10}{fF} to VOUTH and less than \SI{1}{fF} to VOUTL, the H half 3 and \SI{7}{fF}.
This asymmetry between the halves is why the taped-out INL is larger.
\begin{table}[H]\centering\scriptsize
\input{table_par_hier.tex}
\caption{As Table~\ref{tab:par_tapeout}, hierarchical re-layout.}
\label{tab:par_hier}
\end{table}}{}

\subsection{Comparison with the chip}
The chip measures bit weight errors of \num{-7.8}/\num{-5.7}/\num{+6.7}/\num{-8.4}/\num{+3.5}\,\%
(report of the MSO measurement), an INL of about \num{0.5} LSB rms and an ENOB of \num{3.8} bits at low
frequency. The taped-out post-layout netlist has the same sign on all five bits, errors within about
\SI{1}{\percent} of the measured ones on bits 1 and 4 and within \SI{5}{\percent} on the others; it reaches about two thirds of the measured INL and half of the measured
loss of ENOB (from \SchStEnob{} to \TpStEnob{} bits, against \num{3.8}). The layout parasitics of the
DAC are therefore the main cause of the measured nonlinearity. What remains (in the measurement bit 0 is less
reduced, bit 2 larger and bit 3 smaller) is of the size that random mismatch of the unit capacitors, the
references of the board or the parasitic resistance (not extracted) can give.
}{}

% Discussion of the simulation results. Macros from results.tex.
\section{Discussion}

\subsection{Static transfer}
At the conditions of the test board the schematic converter is ideal within the resolution of the
bisection: DNL and INL stay within \numrange{\DnlMin}{\DnlMax} LSB and \numrange{\InlMin}{\InlMax}
LSB, while the transition levels are known to \SI{\Resolution}{mV}, about \num{0.016} LSB. There are
no missing codes, the offset is \num{\OffsetLsb} LSB and the end-point LSB is \SI{\Lsb}{mV} instead of
\SI{\LsbIdeal}{mV} (gain error \SI{\GainErr}{\percent}). A least-squares fit of the transition levels to
binary weights gives
\begin{center}\small
\begin{tabular}{lccccc}
\toprule
Bit & 0 & 1 & 2 & 3 & 4 \\
\midrule
$V_{DD} = \SI{\Vdd}{V}$ (board) & \BitWA & \BitWB & \BitWC & \BitWD & \BitWE \\
\ifdefined\RefVdd
$V_{DD} = \SI{\RefVdd}{V}$ (design testbench) & \RefBitWA & \RefBitWB & \RefBitWC & \RefBitWD & \RefBitWE \\
\fi
Ideal & 1 & 2 & 4 & 8 & 16 \\
\bottomrule
\end{tabular}
\end{center}
\ifdefined\RefVdd
The testbench conditions of the design (\SI{\RefVdd}{V}, LSB \SI{\RefLsb}{mV} against
\SI{\RefLsbIdeal}{mV}) give the same picture: gain error \SI{\RefGainErr}{\percent}, $|$DNL$|$ and
$|$INL$|$ below \RefDnlMax{} and \RefInlMax{} LSB. The supply and the reference levels therefore do not
change the static linearity of the circuit, and the lower supply of the board does not starve the
switches at these reference levels.
\fi

The unit capacitor is a MOS gate capacitance that changes by about
\SI{10}{\percent} over the bias range (\SI{\CapMin}{fF} to \SI{\CapMax}{fF}), yet it does not show up in
the static transfer. All units are the same device with the same bias history, and the comparator
always decides with the two top plates close to each other, so the ratios of the DAC capacitances, which
set the weights, do not depend on the input. The small gain error, the same at both operating points, is
a ratio effect as well (the capacitance of the switches and of the dummy units on the top plates).

\subsection{Consequences for the chip measurement}
The measured chip (report of the MSO measurement in \texttt{RISCVConsole/software/mso}) has bit
weights that differ from binary by \SIrange{-8}{7}{\percent}, an INL of about \num{0.5} LSB rms and a
gain error of about \SI{6}{\percent}. None of this is present in the schematic simulation at the same
supply, references and clock.
\IfFileExists{results_pex.tex}{%
The post-layout simulation of the taped-out layout (Section~\ref{sec:postlayout}) reproduces the sign
of all five weight errors, about two thirds of the INL and half of the loss of ENOB: the main cause is
the wiring capacitance of the DAC plates, which differs from bit to bit (the bridges of bits 1 and 3,
the coupling of the MSB plate to its own top plate and of the other floating plates to the opposite top
plate) and between the two halves of the taped-out layout. The rest can come from}{%
Possible causes, outside of what is simulated here:}
\begin{itemize}
  \item the references of the board, which have no decoupling at the chip and are driven by
        \SI{1}{MHz} op-amps: how far they recover within one bit time depends on the bits already decided;
  \item random mismatch of the unit capacitors (not simulated);
  \item the wiring resistance, which the extraction leaves out.
\end{itemize}
A measurement with the references decoupled at the chip would separate the board from the chip.

\subsection{Dynamic performance}
\IfFileExists{table_sweep.tex}{%
With the ideal differential input the transient simulation gives \SineLowEnob{} bits at
\SI{\SineLowFreq}{kHz} and \SineHighEnob{} bits at \SI{\SineHighFreq}{kHz}, practically the
\SineLowSinadSt~dB that the static transfer alone predicts (SINAD$_{st}$). The sampling network adds
almost no error up to the Nyquist frequency: at \SI{\SineHighFreq}{kHz} the SINAD is
\SineHighSinad~dB and H3 rises from \SI{\SineHtreeLow}{dBc} to about \SI{\SineHtreeHigh}{dBc}, still far
below the quantization error. H2 is a numerical zero: the two halves of the converter are identical and
the inputs exactly opposite, so the even orders cancel. The records at \SI{16.7}{kHz} and \SI{94.9}{kHz}
give identical results, as they must for a converter without dynamic errors: with a coherent record and
a number of cycles prime to the number of samples, both sample the same set of phases of the sine.
\IfFileExists{results_fe.tex}{\input{results_fe.tex}%

With the inputs of the board front-end model (Section~\ref{sec:fe}) the ENOB falls to \FeEnobHigh{} bits
at \SI{\FeFreqHigh}{kHz}. Most of this is the smaller amplitude (\FeAmpHigh{} LSB instead of
\FeAmpLow{}). The common mode that moves with the signal produces H2, but at most
\SI{\FeHtwoMax}{dBc} (at \SI{\FeHtwoMaxFreq}{kHz}): the converter rejects the common-mode ripple well.
The chip measurement shows H2 between \SI{-39}{dBc} and \SI{-24}{dBc} from \SI{300}{kHz} to
\SI{650}{kHz}, much more than this, so the measured H2 is not the response of the ADC to the
common-mode ripple of a linear front end; it is distortion of the front end itself (the op-amps near
their slew rate limit).}{}}{%
The sine records are still running; this subsection is completed with them.}

\subsection{Limits of this simulation}
Schematic netlist for the dynamic part; the post-layout netlists (coupling capacitance, no
resistance) for the static transfer only. Typical corner at \SI{25}{\celsius}, no mismatch and no
noise; the board front end enters only through a linear model of its amplitude and phase.


\begin{thebibliography}{9}
\bibitem{ieee1241} IEEE Std 1241-2010, \emph{IEEE Standard for Terminology and Test Methods for
  Analog-to-Digital Converters}, IEEE, 2011.
\bibitem{ihppdk} IHP, \emph{SG13G2 open source PDK}. \url{https://github.com/IHP-GmbH/IHP-Open-PDK}
\bibitem{ngspice} H.~Vogt, G.~Atkinson, P.~Nenzi, \emph{Ngspice User's Manual}, version 47.
  \url{https://ngspice.sourceforge.io}
\bibitem{kpex} M.~K\"ohler, \emph{KLayout-PEX}, parasitic extraction for KLayout, version 0.3.12.
  \url{https://iic-jku.github.io/klayout-pex-website/}
\bibitem{sboa216} Texas Instruments, \emph{Single-Ended Input to Differential Output Circuit},
  Analog Engineer's Circuit SBOA216B, rev.\ Oct.\ 2024. \url{https://www.ti.com/lit/pdf/sboa216}
\bibitem{lmv358} Texas Instruments, \emph{LMV321, LMV358, LMV324 Low-Voltage Rail-to-Rail Output
  Operational Amplifiers}, datasheet SLOS263Y, rev.\ Aug.\ 2023. \url{https://www.ti.com/lit/gpn/lmv358}
\bibitem{magic} R.~T.~Edwards, \emph{Magic VLSI Layout Tool}, version 8.3.
  \url{http://opencircuitdesign.com/magic/}
\bibitem{netgen} R.~T.~Edwards, \emph{Netgen}, netlist comparison (LVS), version 1.5.
  \url{http://opencircuitdesign.com/netgen/}
\bibitem{klayout} M.~K\"offerlein, \emph{KLayout}, version 0.30. \url{https://www.klayout.de}
\end{thebibliography}


\end{document}
